\documentclass[preview]{standalone}

\usepackage{amsmath}
\usepackage{amssymb}
\usepackage{stellar}
\usepackage{definitions}

\begin{document}

\id{quotient-spaces-sum-scalar-product}
\genpage

\section{Addition}

\begin{snippetdefinition}{quotient-vector-space-addition-definition}{Addition in a quotient space}
    Let \(V\) be a \vectorspace over a \field \(\mathbb{F}\), and let
    \(W\subseteq V\) be a linear subspace. For \(v_1+W,v_2+W\in V/W\),
    define
    \[
        (v_1+W)+(v_2+W)=(v_1+v_2)+W.
    \]
\end{snippetdefinition}

\begin{snippetproposition}{quotient-vector-space-addition-well-defined}{Addition in a quotient space is well-defined}
    Let \(V\) be a \vectorspace over a \field \(\mathbb{F}\), and let
    \(W\subseteq V\) be a linear subspace. The operation
    \[
        (v_1+W)+(v_2+W)=(v_1+v_2)+W
    \]
    does not depend on the chosen representatives of the two cosets.
\end{snippetproposition}

\begin{snippetproof}{quotient-vector-space-addition-well-defined-proof}{quotient-vector-space-addition-well-defined}{Addition in a quotient space is well-defined}
    Suppose that \(v_1'=v_1+w_1\) and \(v_2'=v_2+w_2\), where
    \(w_1,w_2\in W\). Then
    \[
        v_1'+v_2'=(v_1+v_2)+(w_1+w_2).
    \]
    Since \(W\) is a linear subspace, \(w_1+w_2\in W\). Hence
    \[
        (v_1'+v_2')+W=(v_1+v_2)+W,
    \]
    as required.
\end{snippetproof}

\begin{plainhtml}
    Therefore, when adding two cosets, their most convenient representatives may be chosen.
\end{plainhtml}

\begin{snippetproposition}{quotient-vector-space-additive-abelian-group}{Additive group of a quotient space}
    Let \(V\) be a \vectorspace over a \field \(\mathbb{F}\), and let
    \(W\subseteq V\) be a linear subspace. Under the operation
    \[
        (v_1+W)+(v_2+W)=(v_1+v_2)+W,
    \]
    the \set \(V/W\) is an abelian \group.
\end{snippetproposition}

\begin{snippetproof}{quotient-vector-space-additive-abelian-group-proof}{quotient-vector-space-additive-abelian-group}{Additive group of a quotient space}
    The group laws follow from the corresponding laws in \(V\). Indeed,
    for \(v_1,v_2,v_3\in V\),
    \begin{align*}
        ((v_1+W)+(v_2+W))+(v_3+W)
        &=((v_1+v_2)+v_3)+W \\
        &=(v_1+(v_2+v_3))+W \\
        &=(v_1+W)+((v_2+W)+(v_3+W)).
    \end{align*}
    The identity element is \(0_{V/W}=0_V+W=W\), and the additive inverse
    of \(v+W\) is
    \[
        -(v+W)=(-v)+W.
    \]
    Finally,
    \[
        (v_1+W)+(v_2+W)
        =(v_1+v_2)+W
        =(v_2+v_1)+W
        =(v_2+W)+(v_1+W),
    \]
    so the operation is commutative.
\end{snippetproof}

\section{Scalar multiplication}

\begin{snippetdefinition}{quotient-vector-space-scalar-multiplication-definition}{Scalar multiplication in a quotient space}
    Let \(V\) be a \vectorspace over a \field \(\mathbb{F}\), and let
    \(W\subseteq V\) be a linear subspace. For \(\lambda\in\mathbb{F}\) and
    \(v+W\in V/W\), define
    \[
        \lambda(v+W)=\lambda v+W.
    \]
\end{snippetdefinition}

\begin{snippetproposition}{quotient-vector-space-scalar-multiplication-well-defined}{Scalar multiplication in a quotient space is well-defined}
    Let \(V\) be a \vectorspace over a \field \(\mathbb{F}\), and let
    \(W\subseteq V\) be a linear subspace. The operation
    \[
        \lambda(v+W)=\lambda v+W
    \]
    does not depend on the chosen representative of the coset \(v+W\).
\end{snippetproposition}

\begin{snippetproof}{quotient-vector-space-scalar-multiplication-well-defined-proof}{quotient-vector-space-scalar-multiplication-well-defined}{Scalar multiplication in a quotient space is well-defined}
    Suppose that \(v'=v+w\), where \(w\in W\). For every
    \(\lambda\in\mathbb{F}\),
    \[
        \lambda v'=\lambda v+\lambda w.
    \]
    Since \(W\) is a linear subspace, \(\lambda w\in W\), and therefore
    \[
        \lambda v'+W=\lambda v+W.
    \]
\end{snippetproof}

\begin{snippetproposition}{vector-space-quotient-is-vector-space}{A quotient space is a vector space}
    Let \(V\) be a \vectorspace over a \field \(\mathbb{F}\), and let
    \(W\subseteq V\) be a linear subspace. The quotient \(V/W\), with
    \[
        (v_1+W)+(v_2+W)=(v_1+v_2)+W,
        \qquad
        \lambda(v+W)=\lambda v+W,
    \]
    is a \vectorspace over \(\mathbb{F}\).
\end{snippetproposition}

\begin{snippetproof}{vector-space-quotient-is-vector-space-proof}{vector-space-quotient-is-vector-space}{A quotient space is a vector space}
    Addition and scalar multiplication are well-defined by the preceding
    propositions, and \(V/W\) is an abelian \group under addition. It remains
    to verify the scalar multiplication axioms. For
    \(\lambda,\mu\in\mathbb{F}\) and \(v,v_1,v_2\in V\),
    \begin{align*}
        \lambda((v_1+W)+(v_2+W))
        &=\lambda((v_1+v_2)+W) \\
        &=\lambda(v_1+v_2)+W \\
        &=(\lambda v_1+W)+(\lambda v_2+W), \\
        (\lambda+\mu)(v+W)
        &=(\lambda v+\mu v)+W \\
        &=(\lambda v+W)+(\mu v+W), \\
        \lambda(\mu(v+W))
        &=\lambda(\mu v+W) \\
        &=(\lambda\mu)v+W \\
        &=(\lambda\mu)(v+W),
    \end{align*}
    and
    \[
        1_{\mathbb{F}}(v+W)=v+W.
    \]
    Hence all \vectorspace axioms hold.
\end{snippetproof}

\end{document}
